\section{Discussion}
\label{sec:discussion}

\subsection{Key Findings}
\label{sec:key_findings}

Our systematic development cycle demonstrates that **deepening the 3D ResNet architecture** combined with **modern training techniques** (mixup, label smoothing, cosine annealing) yields substantial gains in DaT-SPECT classification, even on a challenging multi-site dataset with significant scanner heterogeneity. The final \vxxv\ model achieves \textbf{OOF LogLoss = 0.2453 (calibrated)} and \textbf{AUROC = 0.9610}, representing an 18.1\% LogLoss reduction and 2.1\% AUROC gain over the strong \vxxiv\ baseline.

\subsection{Why Deeper Architecture Works}
\label{sec:why_deeper}

The transition from 2 to 3 residual blocks + downsampling stages (effectively adding a third hierarchical level) allows the network to capture features at three spatial scales:
\begin{itemize}
    \item Stage 1 (40$^3$): Local texture and edge detectors
    \item Stage 2 (20$^3$): Striatal sub-region patterns (putamen vs caudate)
    \item Stage 3 (10$^3$): Global uptake asymmetry and shape
\end{itemize}

This mirrors the known anatomy: DaT-SPECT signal loss begins in the posterior putamen, spreads to anterior putamen and caudate, with characteristic left-right asymmetry patterns that require multi-scale integration.

\subsection{Scanner Heterogeneity Robustness}
\label{sec:scanner_discussion}

The largest \vxxv\ gains occur on lower-resolution scans (+2.3\% AUROC on 3.90~mm and other). This suggests that:
\begin{enumerate}
    \item The deeper architecture learns resolution-invariant features through mixup augmentation and multi-scale processing.
    \item The coarse-to-fine alignment (rotation search $\pm$5$^\circ$ + translation $\pm$4 voxels) is more effective when the network itself is robust to misalignment.
    \item Per-volume z-score normalization (vs. dataset-level) prevents scanner-specific intensity distributions from biasing the model.
\end{enumerate}

\subsection{Calibration for Clinical Deployment}
\label{sec:cal_discussion}

The 57\% ECE reduction (0.0142 $\to$ 0.0061) via Platt scaling is critical for clinical use. In PD diagnosis, the predicted probability directly informs:
\begin{itemize}
    \item **Risk stratification:** Patients with $p > 0.8$ may be fast-tracked for neurologist referral.
    \item **Longitudinal monitoring:** Probability changes over time track disease progression.
    \item **Trial enrichment:** Well-calibrated probabilities enable statistically valid cohort selection.
\end{itemize}

Uncalibrated models tend toward overconfidence (ECE=0.0142), which could lead to false certainty in borderline cases. Our calibrated probabilities maintain AUROC while providing reliable absolute risk estimates.

\subsection{Computational Efficiency on 4GB VRAM}
\label{sec:compute_discussion}

Training 1.3M-parameter models on 4GB VRAM required several techniques:
\begin{itemize}
    \item Gradient accumulation ($\times$3) $\to$ effective batch 24
    \item Mixed precision (AMP) $\to$ 2$\times$ memory reduction
    \item Per-volume loading (no cached batches) $\to$ minimal CPU RAM
    \item Early stopping at epoch 60--80 $\to$ avoids overfitting and saves time
\end{itemize}

This demonstrates that modern deep 3D CNNs are feasible on consumer hardware with proper engineering, democratizing medical AI research.

\subsection{Limitations}
\label{sec:limitations}

\begin{enumerate}
    \item \textbf{Single tracer:} All scans use $^{123}$I-FP-CIT (DaTSCAN). Generalization to other tracers (e.g., $^{99m}$Tc-TRODAT-1) is untested.
    \item \textbf{Binary classification only:} Does not differentiate PD from atypical parkinsonism (MSA, PSP) which also show dopaminergic deficit.
    \item \textbf{No clinical metadata:} Age, sex, disease duration, and UPDRS scores could improve performance but were not available.
    \item \textbf{Smoke test class shift:} 65\% positive vs 55\% training causes calibration mismatch on small test sets.
    \item \textbf{No external validation:} All results are on the challenge dataset; independent multi-center validation is needed.
    \item \textbf{Ensemble inference time:} 10 models $\times$ 15 TTA views = 150 forward passes ($\sim$8 sec/scan on RTX 3050 Ti). Distillation to a single model would accelerate deployment.
\end{enumerate}

\subsection{Future Work}
\label{sec:future_work}

\begin{itemize}
    \item \textbf{Model distillation:} Compress 10-model ensemble into 1-2 models via knowledge distillation for real-time inference.
    \item \textbf{Multi-task learning:} Joint classification + SBR regression + striatal segmentation for interpretability.
    \item \textbf{Domain adaptation:} Test on external datasets (PPMI, local hospital) with adaptive normalization.
    \item \textbf{Uncertainty quantification:} Monte Carlo dropout or deep ensembles for epistemic uncertainty estimates.
    \item \textbf{Longitudinal modeling:} Incorporate serial scans for progression prediction.
\end{itemize}